\color{unimain}{\subsection{Логика отключения выключателей ВН и ТМП}}
\color{black}

\begin{enumerate}

\item
Устройство формирует выходное воздействие на отключение В1 ВН, В2 ВН и ТМПд по схеме «ИЛИ» в следующих случаях:
\begin{itemize}
\item Срабатывание ДТЗ;
\item Срабатывание вторых ступеней МТЗ, ТЗНП, ТЗОП КО;
\item Срабатывание второй ступени МТЗ НН;
\item Срабатывание МТЗ ОУ;
\item Срабатывание КИ НН;
\item Отключение от внешних защит НН (ЗДЗ, УРОВ ТМП).
\end{itemize}

\item
Устройство формирует выходное воздействие на отключение ТМПо, ТМПр при срабатывании первых ступеней токовых защит КО и первой ступени МТЗ НН.

\item
Алгоритм работы логики отключения В1 ВН приведен на рисунке \ref{pict_lo:LO}. Алгоритмы работы логики отключения В2 ВН и ТМПд выполнены аналогично. Уставки ЛО В1 ВН, В2 ВН и ЛО ТМПд приведены в таблице \ref{app_set:lov1}-\ref{app_set:lotmpd}.

\medskip
\begin{figure}[H]
\centering
\includegraphics[width=1\textwidth,height=1\textheight,keepaspectratio]{\fbpath/11.02 УШР/Логика отключения В/img/LO.pdf}
\captionof{figure}{Функциональная схема логики отключения В1 ВН}\label{pict_lo:LO}
\end{figure}

\item
Алгоритм работы логики отключения ТМПо приведен на рисунке \ref{pict_lo:lo_tmpd}. Алгоритм работы логики отключения ТМПр выполнен аналогично. Уставки ЛО ТМПо, ЛО ТМПр приведены в таблице \ref{app_set:lotmpo}, \ref{app_set:lotmpr}.

\medskip
\begin{figure}[H]
\centering
\includegraphics[width=1\textwidth,height=1\textheight,keepaspectratio]{\fbpath/11.02 УШР/Логика отключения В/img/lo_tmpd.pdf}
\captionof{figure}{Функциональная схема логики отключения ТМПо}\label{pict_lo:lo_tmpd}
\end{figure}


\end{enumerate}